% Class-compatibility preamble retained for a reader excerpt of
%   Sahar Diskin, Anna Geisler, "Perfect Matching in Product Graphs and in their
%   Random Subgraphs", arXiv:2404.14020v4, CC BY 4.0.
%
% The source pins its own publisher class (scrartcl); the delivery wrapper accepts
% only the classes article and amsart, so this file supplies the amsart baseline and
% then copies the author's own declaration block mechanically and unchanged
% (pinned-original-source.tex, source lines 30-57: the author's newtheorem
% declarations and the author's macro declarations).  The author's own package
% list of source lines 4-21 is reused only for the entries the retained text needs
% (enumitem, amsmath, amssymb, amsthm, mathtools, xcolor, geometry) and the
% cross-reference packages already needed by the author's own \Cref/\eqref usage
% (hyperref, cleveref); tikz, caption, subcaption, float, graphicx, wrapfig and
% abstract serve figures and the title block that this excerpt does not carry.
% No mathematical macro, statement, symbol or assumption is introduced, changed
% or removed by this file.
\documentclass[10pt,a4paper]{amsart}
\usepackage[shortlabels]{enumitem}
\usepackage{amsmath}
\usepackage{amssymb}
\usepackage{amsthm}
\usepackage{mathtools}
\usepackage{xcolor}
\usepackage[margin=1in]{geometry}
\usepackage{hyperref}
\usepackage{cleveref}
\hypersetup{
  colorlinks=true,
  linkcolor=blue,
  filecolor=magenta,
  urlcolor=cyan
}

\newtheorem{theorem}{Theorem}
\newtheorem{thm}{Theorem}[section]
\newtheorem{corollary}[thm]{Corollary}
\newtheorem{lemma}[thm]{Lemma}
\newtheorem{proposition}[thm]{Proposition}
\newtheorem{claim}[thm]{Claim}
\newtheorem{remark}[thm]{Remark}
\newtheorem{question}[thm]{Question}

\newcommand{\whp}{whp}
\newcommand{\defi}{\operatorname{def}}
\newcommand{\odd}{\operatorname{odd}}
\newcommand{\iso}{\operatorname{iso}}
\newcommand{\cN}{\mathcal{N}}
\newcommand{\cO}{\mathcal{O}}
\newcommand{\cE}{\mathcal{E}}
\newcommand{\cK}{\mathcal{K}}
\newcommand{\cP}{\mathcal{P}}
\newcommand{\cQ}{\mathcal{Q}}

\newcommand{\footremember}[2]{%
  \footnote{#2}%
  \newcounter{#1}%
  \setcounter{#1}{\value{footnote}}%
}

\newcommand{\ag}[1]{\textcolor{purple}{(Anna: #1)}}
\newcommand{\sd}[1]{\textcolor{red}{(Sahar: #1)}}
\crefname{originallemma}{original lemma}{original lemmas}
\Crefname{originallemma}{Original lemma}{Original lemmas}
\makeatletter
\newlabel{lem:rooted-trees}{{2.3}{source}{}{sourcepointer-1}{}}
\newlabel{lem:rooted-trees@cref}{{[originallemma][0][]2.3}{[1][0][]source}{}{}{}}
\makeatother
\crefname{originalequation}{original equation}{original equations}
\Crefname{originalequation}{Original equation}{Original equations}
\makeatletter
\newlabel{eq:mu-range}{{4.2}{source}{}{sourcepointer-2}{}}
\newlabel{eq:mu-range@cref}{{[originalequation][0][]4.2}{[1][0][]source}{}{}{}}
\makeatother
\title{In the random subgraph of a high-dimensional product graph there is typically no component whose order lies in the interval from two to d to the twentieth}
\author{Sahar Diskin, Anna Geisler}
\date{arxiv:2404.14020v4; CC BY 4.0}
\begin{document}
\maketitle
\noindent\textit{Editorial note.} All delivered mathematical text is the author's own, copied from the pinned author source sources/combinatorics-fl04-vspace-1-5cm-perfect-matching-in-p/pinned-original-source.tex; no statement or proof step was rewritten, reordered or removed.\par\medskip
\noindent\textit{Editorial note.} Preamble provenance: the source pins the publisher class scrartcl, which the delivery wrapper does not accept (supported wrapper classes are article and amsart), so this card uses sources/combinatorics-fl04-vspace-1-5cm-perfect-matching-in-p/retained-preamble.tex, which supplies the amsart baseline and then carries the author's own newtheorem and macro declarations (pinned source lines 30-57) verbatim. No macro was invented and no mathematical symbol was changed.\par\medskip
\noindent\textit{Editorial note.} Bibliography provenance: the source itself carries no thebibliography environment and ends with bibliographystyle abbrv plus bibliography perc (author source lines 1691-1692), so the bibliography is resolved with BibTeX; the database used is the identical copy staged with the same arXiv version under sources/product-graph-matchings-2404.14020v4/perc.bib, whose licence evidence is also status supported for source\_id arxiv:2404.14020v4. No entry was edited and the bibliography style is the author's own abbrv.\par\medskip
\noindent\textit{Editorial note.} The displayed isoperimetric inequality that the counted proof invokes, source label eq:expansion-log, is delivered verbatim as a context span so that the first line of the proof refers to a formula present in the assembled material; it is the author's own line of the author's displayed pair of inequalities, quoted in place.\par\medskip
\noindent\textit{Editorial note.} The rooted-tree bound used in the proof and the displayed range for p are navigation pointers only; their content is declared as named standard prerequisites and they are not reproduced.\par\medskip
\noindent\textit{Allowed standard prerequisites (editorial).}\begin{itemize}
\item The parameters n, d, p and the constant M, the relation d = Theta\_M(log n), and the low-degree expansion estimates of the product graph. Set up in the author's preliminaries and proof outline (source lines 201-240); the hypothesis n(1-p)\textasciicircum{}d <= d\textasciicircum{}\{1/50\} used at the end of the proof is the author's displayed eq:mu-range cited below as a pointer. Supplies the expansion bound e\_G(K,K\textasciicircum{}c) >= k(d - (M-1) log\_M k) applied to a set K of size k, and the parameter relations that turn the probability bound into n\textasciicircum{}\{1-0.8k\}.
\end{itemize}\medskip
\ifdefined\hypertarget\hypertarget{sourcepointer-1}{}\fi
\noindent\textit{Source reference.} Original lemma 2.3: sources/combinatorics-fl04-vspace-1-5cm-perfect-matching-in-p/pinned-original-source.tex, source lines 231--231. Navigation pointer only: the counted proof uses the rooted-tree bound in one sentence to control the number of spanning trees; that bound is declared as a named standard prerequisite and is not reproduced here.\par\medskip
\ifdefined\hypertarget\hypertarget{sourcepointer-2}{}\fi
\noindent\textit{Source reference.} Original equation 4.2: sources/combinatorics-fl04-vspace-1-5cm-perfect-matching-in-p/pinned-original-source.tex, source lines 479--479. Navigation pointer only: the counted proof uses the range of p displayed there in one step of the substitution; the display is delivered verbatim in this lane's card for lem:isolates-far but is not part of this card.\par\medskip
\begin{lemma}\label{lem:component-gap}
Whp, $G_p$ has no component whose order lies in $[2,d^{20}]$.
\end{lemma}

\begin{align}
  e_G(A,A^c)&\ge \frac{e}{M}|A|\ln\frac{n}{|A|},\label{eq:expansion-log}\\
  e_G(A,A^c)&\ge |A|\bigl(d-(M-1)\log_M|A|\bigr).\label{eq:expansion-local}
\end{align}

\begin{proof}
Fix $2\le k\le d^{20}$.
For every set $K\subseteq V$ of size $k$, the expansion from \eqref{eq:expansion-local} gives
\[
  e_G(K,K^c)\ge k\bigl(d-(M-1)\log_Mk\bigr)\ge0.9kd.
\]
If $K$ is a component of $G_p$, it contains a spanning tree all of whose edges are present, while every edge of $E_G(K,K^c)$ is absent.
Thus, using \Cref{lem:rooted-trees}, we obtain
\begin{align*}
  \mathbb{P}(G_p\text{ has a component of order }k)
  \le n(ed)^{k-1}(1-p)^{0.9kd} \le n^{1-0.8k},
\end{align*}
Where the last inequality follows since $2\le k\le d^{20}$,  $n (1-p)^d\le d^{1/50}$ and $d=\Theta_M(\ln n)$.
Summing over $k$ shows that whp there is no component whose order lies in $[2, d^{20}]$.
\end{proof}

\bibliographystyle{plain}
\bibliography{perc}
\end{document}
